\documentclass[a4paper,twoside]{ctexrep} % 使用 CTEX 模板
%\usepackage{template/swunthesis-master-professional} %学硕注释掉这一行
\usepackage{template/swunthesis-master} %专硕注释掉这一行
%===============================================
\usepackage{ctex}
\usepackage{fancyhdr}
\usepackage{float} %强制固定图片位置
\usepackage{enumitem}
\usepackage{bicaption}
\usepackage{graphicx}
\usepackage{caption}
\usepackage{titletoc}
\usepackage{tabularx}
\usepackage{lipsum} % 使用时删除

\title{基于\LaTeX 的西南民族大学（专业）硕士学位论文排版格式}
\entitle{The Typesetting Format of (Professional) Postgraduate Dissertation of Southwest Minzu University based on \LaTeX }

\author{于博洋}
\enauthor{Boyang Yu}

\adviser{蔡英}
\coadviser{XXX} % 校外导师（only专业硕士学位）


\college{电子信息学院}
\encollege{College of Electronics and Information}

\major{人工智能} % only学术硕士学位
\research{电子信息}  %学科门类 or 类别（领域）
\resdirection{计算机视觉} % 研究方向
\grade{2020级}

\findate{2024年x月x日} %完成日期
\date{\today} %定稿日期、签署授权书日期

\usepackage{fancyhdr}
\DeclareCaptionLabelFormat{english}{#1~\textbf{#2}}

%设置目录样式，不建议修改
%%%%%%%%%%%%%%%%%%%%%%%%%%%%%%%%%%%%%%%%%%%%%%%%%%%%%%%%%%%%%%%%%%%
\titlecontents{chapter}[3em]{\zihao{4}\songti}{\contentslabel{2.5em}\hspace{1em}}{}{\titlerule*[8pt]{.}\contentspage}
% 设置节目录样式为小四号宋体
\titlecontents{section}[5em]{\zihao{-4}\bfseries \songti}{\contentslabel{2.5em}}{}{\titlerule*[8pt]{.}\contentspage} 
% 设置小节目录样式为小四号宋体
\titlecontents{subsection}[6.5em]{\zihao{-4}\songti}{\contentslabel{3em}}{}{\titlerule*[8pt]{.}\contentspage}
%%%%%%%%%%%%%%%%%%%%%%%%%%%%%%%%%%%%%%%%%%%%%%%%%%%%%%%%%%%%%%%%%%%
\begin{document}
\makecover

\begin{study}

\end{study}


% 摘要
\begin{abstract}{关键词1；关键词2；关键词3}  %关键词位置
这是一个中文摘要示例。
\end{abstract}

\begin{abstractEng}{Keywords 1; Keywords 2; Keywords 3}
This is an example of Abrtract in English.
\end{abstractEng}


\begin{tocc}

\end{tocc}

\begin{foreword}
当撰写论文的前言时，应该包括一系列内容以引导读者对论文有清晰的认识。首先，前言应该简要介绍研究领域的背景和当前的研究现状，明确指出该研究所处的学术位置和价值所在，以及解决的具体问题或探索的内容。其次，要明确阐述论文的研究目的和意义，即研究的目标、动机和预期成果，让读者了解研究的价值和意义所在。接着，论文前言还应描述研究所采用的方法和技术路线，说明研究的思路和实施过程，确保读者对研究过程有清晰的认识。此外，还要对论文的结构安排进行简要介绍，包括各章节的主要内容和组织结构，以便读者更好地理解论文的整体架构和内容安排。最后，在前言中还可以包括对研究过程中获得支持和帮助的感谢之词，向给予帮助和支持的个人、机构表示诚挚的感谢。综上所述，论文的前言应该清晰明了地阐述研究背景、目的和意义，描述研究方法和技术路线，介绍论文的结构安排，同时表达对支持和帮助的感谢之情，为读者提供良好的阅读导引和理解框架。
\end{foreword}

\zihao{-4}
%按照毕业论文格式，重新设置页码（建议不要删除）
\cfoot{\thepage} % 页码居中显示
\setcounter{page}{1}
\pagestyle{fancy}

% 正文示例，按下方样式引入新的章节
\chapter{绪~~论}
\section{示例章节1}
\zihao{-4}
这是一个示例章节.

\lipsum[1-3] % 正式使用模版时删除
\footnote{这里是脚注内容。}
\section{国内外研究现状}
这是一个示例引用。\textsuperscript{\cite{dongshihai2004}}

\subsection{正文示例}
\zihao{-4}
以下是一个正文内容示例。

\begin{enumerate}
	\item 这是一个有序号排列。
	\item 同上。
\end{enumerate}

\begin{itemize}
	\item 这是一个无序号排列。
	\item 同上。
	
\end{itemize}


\begin{figure}[htbp!]
	\centering
	\includegraphics[scale=0.5]{figures/logo.jpg}
	\caption{这是图片的中文标题}
    \captionsetup{labelformat=english}
    \caption*{\bfseries Figure~1-1~This is a sample of image} %需要手动修改
    \label{fig1}
\end{figure}
\begin{table}[H]
	\centering
	\caption{这是一个表格示例}
    \captionsetup{labelformat=english}
    \caption*{\bfseries Table~1-1~This is a sample of table.}
	\begin{tabularx}{\textwidth}{>{\centering\arraybackslash}X>{\centering\arraybackslash}X>{\centering\arraybackslash}X}
			\toprule
			名称 & 名称 & 名称 \\ \midrule
			内容 & 4, 5, 6 & 123 \\
			内容 & 4, 5, 6 & 123 \\ \bottomrule
	\end{tabularx}
\end{table}
\chapter{示例章节}
\lipsum[3] % 正式使用模版时删除
\footnote{这里是脚注内容2。}

\section{示例Section-2}
\begin{figure}[htbp!]
	\centering
	\includegraphics[scale=0.1]{example/template/title.jpg}
	\caption{这是一个图片示例2}
    \captionsetup{labelformat=english}
    \caption*{\bfseries Figure~2-1~This is a sample of image 2.}%需要手动修改
    \label{fig2}
\end{figure}

\subsection{示例Subsection-2}
\begin{table}[H]
	\centering
	\caption{这是一个表格示例2}
    \captionsetup{labelformat=english}
    \caption*{\bfseries Table~2-1~This is a sample of table 2.} %需要手动修改
	\begin{tabularx}{\textwidth}{>{\centering\arraybackslash}X>{\centering\arraybackslash}X>{\centering\arraybackslash}X}
			\toprule
			名称 & 名称 & 名称 \\ \midrule
			内容 & 4, 5, 6 & 123 \\
			内容 & 4, 5, 6 & 123 \\ \bottomrule
	\end{tabularx}
    \label{tab2}
\end{table}
%结论
\chapter*{结~~论}
这是一个结论示例。


% 附录
\appendix
  \chapter{宏包的使用许可}
  \label{appendix:a}
  \date{2024/03}
  \emph{西南民族大学本科毕业设计(论文)} 宏包由2012级计科1202欧长坤进行开发，本宏包依照 GNU LGPL 协议授权分发的自由软件, 在使用本项目时，您的可以：

\begin{itemize}
    \item 任意下载本项目且无需支付任何费用
    \item 任意将本项目的副本分发给他人使用
    \item 获取和修改本项目的源代码
\end{itemize}

  唯需遵守以下条件：

  \begin{itemize}
      \item 当您将自己对本项目的修改版(即衍生作品)发布时，衍生作品也必须按照GNU LGPL或更严格的协议发布
  \end{itemize}

	\textbf{2024/03 修改更新说明}
	
	基于欧长坤学长开发的\emph{西南民族大学本科毕业设计(论文)} \LaTeX 模板宏包，本次修改版本由电子信息学院2020级 于博洋~进行补充，具体修改内容如下：
	\begin{itemize}
		\item 添加了对于目录部分排版的要求
		\item 精简了整体结构，不再需要从content文件夹导入正文等部分
		\item 整体的排版格式，文档、标题、字号要求进行更新和修改
		\item 重新调整了参考文献部分的排版格式
		\item 使用协议遵循原始版本
	\end{itemize}
	再次感谢欧学长对于SMU-Thesis \LaTeX 模板的开发和贡献,祝愿未来科研与工作一切顺利！也希望后续使用的同学们可以根据需求继续完善Master-Professional \LaTeX 和Master-Research \LaTeX 模板，也希望后续的修改可以继续补充在该附录使用许可中。
	
  \textbf{特别声明}：由于 LaTeX 文档的特殊性，使用此宏包构建的文档（即学位论文）不属于 LGPL 协议生效的一部分，即用户无需在文档中包含本宏包的版权声明和引用说明。如果你希望表示对本项目的支持，可以在学位论文末尾的感谢处保留对项目作者以及后续修改的支持和感谢。

  \cleardoublepage
 % 阅读使用说明许可后，正式使用模板时移除此行代码。

% 参考文献
\bibliographystyle{template/bstutf8}
\clearpage
\phantomsection
\songti
\linespread{1.0}   
\centering
\addcontentsline{toc}{chapter}{参考文献} % 向目录中添加参考文献条目

\bibliography{references/main}

% 致谢
\chapter*{致~~谢}
\addcontentsline{toc}{chapter}{致谢} % 此命令添加不含序号的标题至目录
\raggedright
\setlength{\parindent}{2em} % 设置首行缩进为2em，根据需要调整缩进大小
感谢老师、感谢亲人、感谢朋友。
\end{document}
